\SourceLeaf{32}
\fancyhead[L]{\small OBSERVATIONS OF THE SECOND COMET OF 1813}
\renewcommand{\Article}[1]{\par\medskip\Needspace{4\baselineskip}\pdfbookmark[1]{Art. #1}{cometae-art-#1}\centerline{\textbf{#1.}}\nobreak\smallskip}
\pdfbookmark[0]{Observations of the Second Comet of 1813}{cometae}
\begin{center}
{\large OBSERVATIONS}\par\medskip
{\Large OF THE SECOND COMET OF 1813}\par\medskip
MADE AT THE GÖTTINGEN OBSERVATORY\par\medskip
WITH\par\medskip
SOME NOTES ON THE CALCULATION OF PARABOLIC ORBITS.
\end{center}

I began my own observations of the comet at our observatory on 7 April. My dear colleague, the distinguished \textsc{Harding}, had discovered it on 3 April of this year in the constellation of Poniatowski’s Bull. Here are the determinations I was able to make with a circular micrometer fitted to a ten-foot telescope:
\begin{center}\small\setlength{\tabcolsep}{4pt}
\begin{tabular}{r|r|r|r l}
1813&Gött. mean time&App. R.A.&\multicolumn{2}{c}{App. decl.}\\\hline
Apr. 7&$13^{\mathrm h}\;12^{\mathrm m}\;2^{\mathrm s}$&$271^0\;7'\;19''3$&$5^0\;34'\;36''7$&N.\\
9&$13\;35\;40$&$270\;10\;33.5$&$4\;11\;3.4$&—\\
11&$13\;17\;43$&$269\;1\;19.9$&$2\;33\;0.7$&—\\
14&$13\;7\;36$&$266\;44\;5.5$&$0\;33\;0.8$&S.\\
21&$14\;23\;0$&$256\;39\;19.3$&$12\;57\;56.0$&—
\end{tabular}\end{center}

The distinguished \textsc{Harding} later made the following observations with the mural quadrant:
\begin{center}\small\setlength{\tabcolsep}{4pt}
\begin{tabular}{r|r|r|r l}
Apr. 21&$15^{\mathrm h}\;7^{\mathrm m}\;21^{\mathrm s}$&$256^0\;34'\;19''6$&$13^0\;2'\;26''5$&S.\\
24&$14\;22\;50$&$248\;23\;21$&$21\;45\;2$&\\
25&$14\;4\;21$&$244\;44\;42$&$25\;10\;42$&
\end{tabular}\end{center}

On 24 and 25 April the comet was very conspicuous even to the naked eye. On the following nights the sky was covered with clouds, and soon the comet’s rapid descent southward put an end to the observations.

\SourceLeaf{33}
It seems unnecessary to repeat here the parabolic elements that I had myself derived from the first three observations at the very outset. I entrusted the task of refining them to a highly experienced and skilled calculator, Dr. \textsc{Gerling}, to whom we owe the following corrected elements. They were fitted as closely as possible to all the observations, both our own and those sent by the distinguished \textsc{Olbers}.
\begin{center}\fontsize{9}{11}\selectfont\setlength{\tabcolsep}{1.5pt}
\begin{tabular}{@{}l l r@{}}
Logarithm of the perihelion distance&. . . . . . . . . .&$0.0849212$\\
\multicolumn{3}{@{}p{\linewidth}@{}}{Time of passage through perihelion, for the Göttingen meridian}\\
\multicolumn{3}{r@{}}{1813 May $19.44507$}\\
Longitude of perihelion&. . . . . . . . . . . . . .&$197^0\;43'\;7''7$\\
Longitude of the ascending node&. . . . , . . . . . .&$42\;40\;15.2$\\
Orbital inclination&. . . . . . . . . . . . . . .&$81\;2\;11.8$\\
\multicolumn{3}{@{}l}{\textit{Retrograde} motion.}
\end{tabular}\end{center}

\Needspace{11\baselineskip}
The distinguished \textsc{Olbers} made the following observations:
\begin{center}\small\setlength{\tabcolsep}{4pt}
\begin{tabular}{r|r|r|r l}
1813&Bremen mean time&App. R.A.&\multicolumn{2}{c}{App. decl.}\\\hline
Apr. 14&$13^{\mathrm h}\;31^{\mathrm m}\;4^{\mathrm s}$&$266^0\;42'\;51''2$&$0^0\;34'\;22''8$&S.\\
15&$12\;14\;29$&$265\;48\;47.9$&$1\;46\;4.5$&\\
19&$11\;38\;0$&$260\;40\;39.1$&$8\;15\;23.7$&\\
21&$12\;0\;35$&$256\;51\;59.3$&$12\;42\;54.3$&\\
24&$11\;58\;38$&$248\;43\;57.7$&$21\;25\;9.8$&\\
25&$11\;41\;30$&$245\;8\;18.0$&$24\;49\;2.4$&\\
&$12\;5\;38$&$245\;4\;3.0$&$24\;54\;16.4$&
\end{tabular}\end{center}

There was also an observation by the distinguished \textsc{Bouvard}, made at the Imperial Observatory in Paris:
\begin{center}\small\setlength{\tabcolsep}{4pt}
\begin{tabular}{r|r|r|r l}
1813&Paris mean time&App. R.A.&\multicolumn{2}{c}{App. decl.}\\\hline
Apr. 13&$16^{\mathrm h}\;22^{\mathrm m}\;2^{\mathrm s}$&$267^0\;27'\;18''$&$0^0\;24'\;46''$&N.
\end{tabular}\end{center}

The distinguished \textsc{Gerling}’s calculations show how closely the elements given above agree with each of these observations:
\SourceLeaf{34}
\begin{center}\small\setlength{\tabcolsep}{6pt}
\begin{tabular}{r|r|r|l}
&\multicolumn{2}{c|}{Difference}&\\
&R.A.&Decl.&Observer\\\hline
Apr. 7&$+3''8$&$+8''5$&Gauss\\
9&$+2.0$&$+34.3$&Gauss\\
11&$-5.3$&$-17.7$&Gauss\\
13&$-1.6$&$-0.4$&Bouvard\\
14&$-7.4$&$-28.6$&Gauss\\
&$+2.7$&$-8.4$&Olbers\\
15&$-0.9$&$+28.7$&Olbers\\
19&$-25.1$&$+103.9$&Olbers\\
21&$-56.6$&$-59.2$&Olbers\\
&$-30.1$&$-5.2$&Gauss\\
&$-22.8$&$-24.1$&Harding\\
24&$-45.2$&$-41.4$&Olbers\\
&$+0.4$&$-11.6$&Harding\\
25&$-23.3$&$-67.8$&Olbers\\
&$-9.4$&$-27.4$&Olbers\\
&$+9.7$&$+1.4$&Harding
\end{tabular}\end{center}

In this comparison, proper allowance was made for both aberration and parallax.
\begin{center}\rule{.25\linewidth}{.3pt}\end{center}

Let me add a few shortcuts in calculation that I have often found useful in the first determination of a parabolic orbit by the distinguished \textsc{Olbers}’s method. They seem to make this already efficient method even shorter or better suited to numerical calculation. They concern the calculation of the radius vectors and, above all, the chord between the first and last positions. For this purpose the distinguished \textsc{Olbers} uses expressions of the form $\BareRoot{(f+g\rho+h\rho\rho)}$. He determines the coefficients $f$, $g$, $h$ by fairly simple operations, but in most cases these cannot give sufficient precision unless they are carried out with the larger tables of logarithms, calculated to seven or at least six decimal places. I have substituted other expressions that seem somewhat easier to calculate in their own right and have the further advantage that five decimal places, and hence the smaller tables of logarithms, suffice for every operation. The essentials are as follows. Let
\SourceLeaf{35}
\begin{description}
\item[$\odot$, $\odot'$, $\odot''$] the longitudes of the Sun at the first, second and third observations
\item[$R$, $R'$, $R''$] the distances of the Sun from the Earth
\item[$\alpha$, $\alpha'$, $\alpha''$] the geocentric longitudes, and
\item[$\ObservedBeta$, $\ObservedBeta'$, $\ObservedBeta''$] the geocentric latitudes of the comet
\item[$r$, $r'$, $r''$] its distances from the Sun
\item[$\rho$, $\rho'$, $\rho''$] its projected distances from the Earth
\item[$t$, $t'$, $t''$] the times of the observations
\item[$k$] the chord between the first and third positions of the comet
\item[$M=\frac{\rho''}{\rho}$]
\end{description}

With these definitions it is easy to see that
\[
\begin{aligned}{}
[1]\quad r&=\BareRoot{[(\rho\cos\alpha-R\cos\odot)^2}\\
&\qquad +(\rho\sin\alpha-R\sin\odot)^2+\rho\rho\operatorname{tang}\ObservedBeta^2]\\
[2]\quad r''&=\BareRoot{[(M\rho\cos\alpha''-R''\cos\odot'')^2}\\
&\qquad +(M\rho\sin\alpha''-R''\sin\odot'')^2\\
&\qquad +MM\rho\rho\operatorname{tang}{\ObservedBeta''}^2]\\
[3]\quad k&=\BareRoot{[(M\rho\cos\alpha''-\rho\cos\alpha-R''\cos\odot''+R\cos\odot)^2}\\
&\qquad +(M\rho\sin\alpha''-\rho\sin\alpha-R''\sin\odot''+R\sin\odot)^2\\
&\qquad +(M\rho\operatorname{tang}\ObservedBeta''-\rho\operatorname{tang}\ObservedBeta)^2]
\end{aligned}
\]
Equations 1 and 2 take the following form
\[
\begin{aligned}
r&=\BareRoot{\left(\frac{\rho\rho}{\cos\ObservedBeta^2}-2\rho R\cos(\alpha-\odot'')+RR\right)}\\
r''&=\BareRoot{\left(\frac{MM\rho\rho}{\cos{\ObservedBeta''}^2}-2M\rho R''\cos(\alpha''-\odot'')+R''R''\right)}
\end{aligned}
\]
Thus, if we set
\[
\begin{aligned}
\cos\ObservedBeta\cos(\alpha-\odot)&=\cos\psi, & R\sin\psi&=B\\
\cos\ObservedBeta''\cos(\alpha''-\odot'')&=\cos\psi'', & R''\sin\psi''&=B''
\end{aligned}
\]
we obtain
\[
\begin{aligned}
r&=\BareRoot{\left[\left(\frac{\rho}{\cos\ObservedBeta}-R\cos\psi\right)^2+BB\right]}\\
r''&=\BareRoot{\left[\left(\frac{M\rho}{\cos\ObservedBeta''}-R''\cos\psi''\right)^2+B''B''\right]}
\end{aligned}
\]

Further, if we determine the five auxiliary quantities $g$, $G$, $h$, $H$, $\zeta$ so that
\SourceLeaf{36}
\[
\begin{aligned}
R''\cos\odot''-R\cos\odot&=g\cos G\\
R''\sin\odot''-R\sin\odot&=g\sin G\\
M\cos\alpha''-\cos\alpha&=h\cos\zeta\cos H\\
M\sin\alpha''-\sin\alpha&=h\cos\zeta\sin H\\
M\operatorname{tang}\ObservedBeta''-\operatorname{tang}\ObservedBeta&=h\sin\zeta
\end{aligned}
\]
formula 3 takes the following form
\[
\begin{aligned}
k&=\BareRoot{[(\rho h\cos\zeta\cos H-g\cos G)^2}\\
&\qquad +(\rho h\cos\zeta\sin H-g\sin G)^2+\rho\rho hh\sin\zeta^2]\\
&=\BareRoot{(\rho\rho hh-2\rho hg\cos\zeta\cos(G-H)+gg)}
\end{aligned}
\]
So if we set
\[
\cos\zeta\cos(G-H)=\cos\varphi,\quad g\sin\varphi=A
\]
we have
\[
k=\BareRoot{[(\rho h-g\cos\varphi)^2+AA]}
\]
or, if we also set $\rho h-g\cos\varphi=u$,
\[
k=\BareRoot{(uu+AA)}
\]

We think many readers will welcome it if we give all the operations involved in these transformations again, in their proper order, and also add every remaining operation, so that everything needed for the first calculation of a parabolic orbit can be found together here. We will also illustrate these instructions with numerical values for our comet. We therefore choose our observations of 7, 14 and 21 April. Reducing them gives the following \textit{data}:
\[
\begin{array}{rcrrr@{\qquad}rcrrr}
t&=&\multicolumn{3}{l}{7.55002}&&&&&\\
t'&=&\multicolumn{3}{l}{14.54694}&&&&&\\
t''&=&\multicolumn{3}{l}{21.59931}&&&&&\\
\alpha&=&271^0&16'&38''&\ObservedBeta&=&+29^0&2'&0''\\
\alpha'&=&266&27&22&\ObservedBeta'&=&+22&52&18\\
\alpha''&=&256&48&8&\ObservedBeta''&=&+9&53&12\\
\odot&=&17&47&41&\log R&=&\multicolumn{3}{l}{0.00091}\\
\odot'&=&24&38&45&\log R'&=&\multicolumn{3}{l}{0.00175}\\
\odot''&=&31&31&25&\log R''&=&\multicolumn{3}{l}{0.00260}
\end{array}
\]

\SourceLeaf{37}
I. The \textit{first} operation is to determine an approximate value of $M$, given by the formula
\[
M=\frac{t''-t'}{t'-t}.\frac{\operatorname{tang}\ObservedBeta'\sin(\alpha-\odot')-\operatorname{tang}\ObservedBeta\sin(\alpha'-\odot')}{\operatorname{tang}\ObservedBeta''\sin(\alpha'-\odot')-\operatorname{tang}\ObservedBeta'\sin(\alpha''-\odot')}
\]
In our example we find $\log M=9.75799$.

II. Next, determine the quantities $g$, $G$, $h$, $H$, $\zeta$ by the following formulas. These are evidently equivalent to those given above and are easier to use in calculation.
\[
\begin{aligned}
R''\cos(\odot''-\odot)-R&=g\cos(G-\odot)\\
R''\sin(\odot''-\odot)&=g\sin(G-\odot)\\
M-\cos(\alpha''-\alpha)&=h\cos\zeta\cos(H-\alpha'')\\
\sin(\alpha''-\alpha)&=h\cos\zeta\sin(H-\alpha'')\\
M\operatorname{tang}\ObservedBeta''-\operatorname{tang}\ObservedBeta&=h\sin\zeta
\end{aligned}
\]

In our example this gives
\[
\begin{array}{rcrrr}
G&=&113^0&43'&57''\\
\log g&=&\multicolumn{3}{l}{9.38029}\\
H&=&109^0&5'&49''\\
\zeta&=&44&13&9\\
\log h&=&\multicolumn{3}{l}{9.81477}
\end{array}
\]

III. Next we set
\[
\begin{aligned}
\cos\zeta\cos(G-H)&=\cos\varphi\\
\cos\ObservedBeta\cos(\alpha-\odot)&=\cos\psi\\
\cos\ObservedBeta''\cos(\alpha''-\odot'')&=\cos\psi''\\[5pt]
g\sin\varphi&=A\\
R\sin\psi&=B\\
R''\sin\psi''&=B''
\end{aligned}
\]
If the cosines of the angles $\varphi$, $\psi$, $\psi''$ happen to differ little from unity in this calculation, it is advisable to use six or even seven places. There is no need to express the angles $\varphi$, $\psi$, $\psi''$ themselves in degrees, minutes and seconds, since it is enough to pass directly from the logarithms of the cosines to those of the sines.

\SourceLeaf{38}
In our example this gives
\[
\begin{aligned}
\log A&=9.22527\\
\log B&=9.98706\\
\log B''&=9.86038
\end{aligned}
\]

IV. Finally, set
\[
\begin{aligned}
h\cos\ObservedBeta&=b\\
\frac{h\cos\ObservedBeta''}{M}&=b''\\
g\cos\varphi-bR\cos\psi&=c\\
g\cos\varphi-b''R''\cos\psi''&=c''
\end{aligned}
\]
In our example
\[
\begin{aligned}
\log b&=9.75645\\
\log b''&=0.05028\\
c&=+0.31365\\
c''&=+0.95443
\end{aligned}
\]

V. With these preparations, the radius vectors $r$, $r''$ and the chord $k$ depend on the unknown $u$ as follows
\[
\begin{aligned}
r&=\BareRoot{\left[\left(\frac{u+c}{b}\right)^2+BB\right]}\\
r''&=\BareRoot{\left[\left(\frac{u+c''}{b''}\right)^2+B''B''\right]}\\
k&=\BareRoot{(uu+AA)}
\end{aligned}
\]
The unknown must be determined by trial so as to satisfy the equation
\[
(r+r''+k)^{\frac32}-(r+r''-k)^{\frac32}=\frac{t''-t}{m}
\]
Here $m$ denotes a time of $9.6887401$ days, and
\[
\log m=0.9862673
\]
The quantity $(r+r''-k)^{\frac32}$ would have to be preceded by the sign $+$ if the heliocentric motion of the comet during the time $t''-t$ exceeded an angle of $180^0$. This case cannot occur under the assumptions on which this first determination of the orbit rests. It scarcely needs saying that in calculating $r$ numerically one should introduce an auxiliary angle $\theta$ such that
\SourceLeaf{39}
\[
\frac{bB}{u+c}=\operatorname{tang}\theta
\]
This gives $r=\frac{B}{\sin\theta}$, and likewise for $r''$ and $k$. Anyone will also readily see that in all these operations our table for finding the logarithms of sums and differences directly can be used to great advantage.

In our example we have $\log\frac{t''-t}{m}=0.16139$, and a few trials give
\[
u=0.24388
\]

VI. Once $u$ has been found, we have
\[
\rho=\frac{u+g\cos\varphi}{h},\quad\rho''=M\rho
\]
\begin{center}(in our example $\log\rho=9.80364$, $\log\rho''=9.56163$).\end{center}

The remaining operations are familiar enough, but for completeness it seemed worth adding the other formulas we generally use. Let
\begin{description}
\item[$\lambda$, $\lambda''$] the heliocentric longitudes of the comet at the first and third observations
\item[$\ApproxBeta$, $\ApproxBeta''$] the heliocentric latitudes
\item[$\upsilon$, $\upsilon''$] the longitudes in the orbit
\item[$\AscendingNode$] the longitude of the ascending node
\item[$i$] the orbital inclination, taken between $0$ and $90^0$ if we distinguish direct and retrograde motion in the usual way
\item[$\Pi$] the longitude of perihelion
\item[$T$] the time of passage through perihelion
\item[$q$] the perihelion distance.
\end{description}

VII. The heliocentric positions are found by the formulas
\[
\begin{aligned}
\rho\cos(\alpha-\odot)-R&=r\cos\ApproxBeta\cos(\lambda-\odot)\\
\rho\sin(\alpha-\odot)&=r\cos\ApproxBeta\sin(\lambda-\odot)\\
\rho\operatorname{tang}\ObservedBeta&=r\sin\ApproxBeta\\
\rho''\cos(\alpha''-\odot'')-R''&=r''\cos\ApproxBeta''\cos(\lambda''-\odot'')\\
\rho''\sin(\alpha''-\odot'')&=r''\cos\ApproxBeta''\sin(\lambda''-\odot'')\\
\rho''\operatorname{tang}\ObservedBeta''&=r''\sin\ApproxBeta''
\end{aligned}
\]

\SourceLeaf{40}
The agreement of the values for the radius vectors $r$, $r''$ obtained here with those previously derived from $u$ will serve as a check on the calculation. The motion is direct or retrograde according as $\lambda''$ is greater or less than $\lambda$.

In our example we find
\[
\begin{aligned}
\lambda&=225^0\;4'\;22'',&\ApproxBeta&=+14^0\;51'\;39'',&\log r&=0.13896\\
\lambda''&=223\;6\;55,&\ApproxBeta''&=+2\;49\;28,&\log r''&=0.11068
\end{aligned}
\]

The comet’s motion is therefore retrograde.

VIII. To find the longitude of the node and the inclination, use the formulas:
\[
\begin{aligned}
\pm\operatorname{tang}\ApproxBeta&=\operatorname{tang}i.\sin(\lambda-\AscendingNode)\\
\pm\frac{\operatorname{tang}\ApproxBeta''-\operatorname{tang}\ApproxBeta.\cos(\lambda''-\lambda)}{\sin(\lambda''-\lambda)}&=\operatorname{tang}i.\cos(\lambda-\AscendingNode)
\end{aligned}
\]
where the upper signs refer to direct motion and the lower signs to retrograde motion. The longitudes in the orbit are then obtained from the formulas
\[
\begin{aligned}
\frac{\operatorname{tang}(\lambda-\AscendingNode)}{\cos i\cdot}&=\operatorname{tang}(\upsilon-\AscendingNode)\\
\frac{\operatorname{tang}(\lambda''-\AscendingNode)}{\cos i}&=\operatorname{tang}(\upsilon''-\AscendingNode)
\end{aligned}
\]
where $\upsilon-\AscendingNode$, $\upsilon''-\AscendingNode$ must be taken in the same quadrants as $\lambda-\AscendingNode$, $\lambda''-\AscendingNode$, respectively.

For our comet we find
\[
\begin{array}{rcrrr}
\AscendingNode&=&42^0&40'&8''\\
i&=&81&1&3\\
\upsilon&=&237&43&7\\
\upsilon''&=&225&31&32
\end{array}
\]

IX. The longitude of perihelion and the perihelion distance are given by the formulas
\[
\begin{gathered}
\frac{1}{\BareRoot r}=\frac{1}{\BareRoot q}.\cos\tfrac12(\upsilon-\Pi)\\
\frac{\operatorname{cotg}\tfrac12(\upsilon''-\upsilon)}{\BareRoot r}-\frac{1}{\sin\tfrac12(\upsilon''-\upsilon).\BareRoot{r''}}=\frac{1}{\BareRoot q}.\sin\tfrac12(\upsilon-\Pi)
\end{gathered}
\]
For our comet, $\Pi=197^0\;37'\;51''$, $\log q=0.08469$.

\SourceLeaf{41}
X. Finally, take from \textsc{Barker}’s table the mean motions corresponding to the true anomalies $\upsilon-\Pi$, $\upsilon''-\Pi$ or $\Pi-\upsilon$, $\Pi-\upsilon''$, and call them $M$, $M''$. Then
\[
T=t\mp Mnq^{\frac32}=t''\mp M''nq^{\frac32}
\]
where the upper signs apply if $\upsilon>\Pi$, $\upsilon''>\Pi$ in direct motion, or if $\upsilon<\Pi$, $\upsilon''<\Pi$ in retrograde motion. The lower signs apply in the opposite cases. The quantity $n$ is constant, with logarithm $=0.0398723$. Agreement of the two values of $T$ provides a second check on the calculation.

In our example we find
\[
\begin{aligned}
T&=49.518\\
T&=49.517
\end{aligned}
\]
so that May $19.5175$ can be taken as the time of passage through perihelion.

If the geocentric position at the time of the middle observation is calculated from these elements, the result is a longitude of $266^0\;27'\;15''$ and a latitude of $22^0\;52'\;18''$ north. The longitude differs from the observed value by $7''$, and the latitude agrees exactly.
\par\medskip\begin{center}\rule{.25\linewidth}{.3pt}\end{center}
